\subsection{Pseudorandom Number Generators}

Cryptography needs a lot of random to be secure. 
Computers supply a little bit through core temperature, 
power consumption, mouse movement, and it would be nice 
to expand those random bits into many more random bits. 

Specifically, given $n$ uniformly random bits 
$x = x_1|x_2|\dots|x_n$, we want a deterministic
polynomial-time algorithm $G$ such that $G(x)$
outputs $l(n)$ bits $y=y_1|y_2|\dots|y_{l(n)}$ 
with $l(n) > n$ and $y$ looks "as good as" a truly random string. 
$G: \{0,1\}^n \rightarrow \{0,1\}^{l(n)}$ is called a pseudorandom
generator (PRG). "As good as" means that the string 
will pass all statistical tests that a truly random string 
would pass. You can never generate more information from 
nothing, so the idea of "as good as" means that no efficient 
computer can tell $G(x)$ and a truly random string apart. 
The distributions
\begin{equation}
    \left\{x \leftarrow \left\{0,1\right\}^{l(n)}: G(x) \right\}
\end{equation}
and
\begin{equation}
    \left\{r \leftarrow\{0,1\}^{l(n)}: r\right\}
\end{equation}
are "computationally indistinguishable". 

"Computationally indistinguishable" means that no PPT 
distinguisher algorithm $D$ can tell them apart. 

A deterministic algorithm $G$ is called a pseudorandom
generator if:
\begin{itemize}
    \item $G$ can be computed in polynomial time. 
    \item $|G(x)| > |x|$
    \item $\left\{x \leftarrow \{0,1\}^n : G(x) \right\} \approx_c \left\{U_{l(n)}\right\}$
    where $l(n) = |G(0^n)|$
\end{itemize}
$\approx_c$ denotes computational indistinguishability. 
$U_{l(n)}$ denotes uniform distribution over $\left\{0,1\right\}^{l(n)}$.
$|G(x) = l(n)|$ is the expansion factor of $G$. 